% ID: SEL-D0072-OB-E01
% ESTADO: cotejada. Segunda pasada visual y borrador compilado revisados.
% Fuente: Selectividad/sel_2009_matematicas_aplicadas/mat apli CC SS06.pdf
% SHA-256: b371bdd6ba8b176301371f25945c7fb030327ff4501a6a02aa6f109a3a4f53e1
% Página: 2; opción B; ejercicio original 1.

\begin{enumerate}[label=\alph*)]
\item \textbf{(1.25 puntos)} Plantee, sin resolver, el siguiente problema de programación lineal:

``Una empresa fabrica camisas de dos tipos, A y B. El beneficio que obtiene es de 8 euros por cada camisa que fabrica del tipo A, y de 6 euros por cada una del tipo B. La empresa puede fabricar, como máximo, 100000 camisas, y las del tipo B han de suponer, al menos, el $60\%$ del total. \textquestiondown{}Cuántas camisas debe fabricar de cada tipo para obtener el máximo beneficio?''
\item \textbf{(1.75 puntos)} Represente la región definida por las inecuaciones:
\[
y\leq x,\quad y+2x\leq6,\quad x\leq4y+3.
\]
Calcule el máximo de $F(x,y)=y+2x$ en la región anterior e indique dónde se alcanza.
\end{enumerate}
